\section{Discussion}
\label{sec:discussion}

\subsection{Reading the Results through Turing, Popper and Occam}
\label{sec:disc-mendel}

Mendel's three principles (Section~\ref{sec:bg-cww}) give a compact reading of the results.
In the sense of \emph{Turing}, an LLM codebook is adequate if its word meanings cannot be told apart from those of
the people it serves.
In the middle of the scale, LLM meanings overlap moderately with those of respondents; for several extremes and negations
they clearly fail this test:
a reader who receives \word{Will Happen} or \word{Improbable} from the LLMs would reconstruct a number far from what
respondents mean (Section~\ref{sec:res-rq2}).
The upward shift of negative phrases matches the bias that a recent comparison of many LLMs reports for
\word{unlikely} and \word{improbable} \citep{petridis2026unlikely}, which suggests that it is not specific to our four
models.

In the sense of \emph{Popper}, a T1 word model is an incorrect theory of a word because it is completely certain once
its parameters are fixed.
Several LLM word models are more certain still: in \WtwoCollapsesPOne{} of \WtwoCellsPOne{} cells a model gave one fixed
number, even at temperature \WTemperature{}, and its word model is a point.
Such a model cannot represent that a word means different things to different people, and the IT2 envelope over models
is the minimal correction: it keeps the disagreement between models as a footprint instead of averaging it away.

In the sense of \emph{Occam}, the simpler of two models should be chosen when both give the same results.
Excluding inattentive respondents gives almost the same citizen codebook (Section~\ref{sec:res-rq1}), so the simpler
codebook without the exclusion suffices.
Occam does not decide between rules or codebooks that give different results.
The primary nearest-centre fallback and the post-hoc neighbour fallback lead to different flip rates
(\WthreeActionFlipRateLlmVsCitizen{} against \WthreePosthocBSevereNeighbourFallbackActionFlipRateLlmVsCitizen{}), so the
choice between them is a design decision about how a service should act where no word fits, and the same holds for the
choice between the citizen and the LLM codebook.

\subsection{Why Extremes Diverge More than Hedges}
\label{sec:disc-shape}

H2a expected the opposite of what we found.
We read the reversal as an effect of the shape of the sets rather than of their position; the study does not test
this explanation directly.
Respondents use extreme phrases narrowly: their word models are thin shoulders close to $0$ or $100$, sometimes nearly
points.
LLMs also answer extreme phrases with one or a few round numbers, but not always the same ones; a point set that is a
few percentage points away from a narrow citizen set has almost no overlap, so the Jaccard similarity drops sharply.
Hedge phrases are wide for respondents, and a narrower LLM set placed anywhere inside them still overlaps substantially.
A small error in position therefore costs more overlap at the extremes than in the middle, and narrowness alone (H2b)
does not make LLM meanings worse for hedges.
For a service, this matters because the extreme words are exactly the ones that trigger the strongest actions.

\subsection{Implications: Codebook Governance}
\label{sec:disc-governance}

If a public service lets an AI system communicate urgency in words, the meanings of those words become part of the
service.
Our results show that the meanings an LLM brings are not those of the people it serves, that they can have gaps (the LLM
codebook covers no level word between \WthreeLlmGapLow{} and \WthreeLlmGapHigh{}), and that the difference can change
whether a helpline is suggested.
We therefore argue for \emph{codebook governance}: a service should define its risk words with the people it serves,
publish the resulting codebook, test the system's use of the words against it, and repeat the test when the model
changes.
In a cognitive city, such a codebook is public data in the sense of Spycher and Portmann: a fuzzy, linguistic
category system that should be maintained openly rather than fixed crisply by the system builder
\citep{spycher2026crispness}.
The \Perc{} makes this practical, because the codebook is a separate, exchangeable component, and the word-meaning lab
shows how the result can be explained to non-experts.
On this data set, the word mismatch changes whether a helpline is suggested, never whether a human is involved; this is
reassuring for the escalation to people that both 143 and 113 require, but it also shows where a mismatch is most likely
to go unnoticed.

\subsection{Limitations}
\label{sec:disc-limitations}

\begin{itemize}
\item \textbf{Point data, not intervals.} CAPphrase records one number per respondent and phrase. Our FOU therefore
captures disagreement between demographic groups, not the uncertainty of individuals about a word
(Section~\ref{sec:concept-encoder}). Interval data would allow the IA or EIA.
\item \textbf{Language and respondents.} The phrases are English, and most respondents live in the United Kingdom or the
United States and hold a university degree. The codebook is not that of the users of a Swiss crisis line, who mostly
speak German, French or Italian.
\item \textbf{Models and sampling.} We tested four small open-weight models locally, with two prompts and one
temperature. Larger or commercial models, other prompts and other temperatures may mean different things, and model
updates can change the codebook.
\item \textbf{Decision design.} The mapping of crisis levels to the four survey words, the rule table and the fallback
are design choices; the rule table is not clinically validated. The crisis degree is the share of LLM labellers, not a
validated risk score, and it takes only a few distinct values.
\item \textbf{Messages.} The BHH messages come from public mental-health datasets, not from crisis-line conversations,
and the reply words come from LLM judges.
\item \textbf{Inference.} All tests are exploratory and uncorrected for multiplicity, the fallback sensitivity analyses
are post hoc, and the number of flips is small.
\end{itemize}

\subsection{Ethics}
\label{sec:disc-ethics}

The study uses public data only: aggregate survey data under an open licence, help-seeking messages from a research
data set used without redistribution and without quoting any message, and answers that we generated from local models.
No person was contacted, and all reported results are aggregates.
The prototype is a research tool; it states on every page that its rules are not clinically validated, and it must not
be used to make decisions about real people in crisis.
A system that suggests a helpline less often to people whose message is ambiguous can harm them; we see the main
ethical contribution of the study in making such effects measurable before deployment.
